\section*{Appendix}
\label{app:contents}
\providecommand{\ours}{\arbor}
\startcontents[sections]
\setupappendixtoc
{\hypersetup{linkcolor=AppendixTOCSection}
\printcontents[sections]{l}{1}{\setcounter{tocdepth}{3}}}

\section{Limitations and Future Work}
\label{app:limitations-future}

Although the empirical results demonstrate the promise of \arbor{} for
autonomous research, this study has several limitations. We discuss these
limitations below and outline the corresponding directions for future work.

\paragraph{Evaluation scope.}
Our experiments are an initial probe of autonomous research rather than a
complete benchmark for scientific discovery. The current AO task suite covers
model training, harness engineering, and data synthesis, but it does not yet
span the full diversity of research problems. Within AI, future tasks should
include settings such as low-level kernel optimization, pretraining data-mixture
design, and more open-ended system design. Beyond AI, domains such as biology,
mathematics, and physics require benchmarks where valuable hypotheses are harder
to specify and validate. A broader suite should therefore evaluate not only
metric improvement, but also whether the generated ideas are scientifically
meaningful, reproducible, and transferable.

\paragraph{Objective design.}
The present AO interface mainly optimizes a fixed scalar objective defined by a
task-specific evaluator. This is useful for controlled experiments, but it is a
simplification of real research. Scientific objectives are often multi-dimensional:
performance, resource use, robustness, interpretability, novelty, and safety may
all matter, and improving one can hurt another. Future AO systems should support
multi-objective search, explicit constraints, Pareto-style comparison, and
adaptive scheduling between competing criteria. This would also reduce the risk
that an agent overfits to a narrow benchmark metric while missing the broader
research goal.

\paragraph{Idea generation.}
We observe that agents can read evaluation feedback carefully and propose useful
local refinements, but their research ability remains far from that of expert
human researchers. In difficult tasks, they may fail to identify a genuinely new
mechanism, abandon a promising direction after early failures, or reverse-engineer
solutions from observed scores instead of reasoning from first principles. A more
fine-grained study of agent idea formation is therefore needed. Promising
directions include better uncertainty tracking, explicit reuse of negative
evidence, mechanisms for revisiting suspended branches, and training or prompting
methods that encourage causal and first-principles hypotheses rather than only
result-driven fixes.

\paragraph{Cost and infrastructure.}
Long-horizon autonomous research is limited not only by idea quality, but also
by systems engineering. In our runs, performance and efficiency depend on details
such as prompt caching, evaluator scheduling, isolated environment startup,
parallel worktree execution, and the reliability of inter-agent coordination.
Large numbers of model calls, evaluator calls, and artifact handoffs can make a
successful search expensive even when each individual step is simple. Future work
should develop cost-aware tree policies, adaptive evaluator allocation, stronger
caching and checkpointing, and more robust execution infrastructure so that AO
systems can scale without turning search breadth into uncontrolled compute cost.

\paragraph{Model capability.}
Finally, \arbor{} inherits the strengths and weaknesses of the underlying LLMs.
Current models are often capable of coding, summarizing results, and making
plausible local hypotheses, but they can still struggle with deep domain
knowledge, long chains of causal reasoning, and genuinely creative problem
reformulation. Stronger foundation models will likely improve AO directly, but
model scaling alone may not be sufficient. Future systems should combine LLMs
with domain knowledge bases, specialized tools, simulators, formal checkers, and
training signals targeted at scientific hypothesis generation. In this sense,
\arbor{} provides a structure for accumulating and testing ideas, while the
quality of those ideas remains an important frontier.

\section{Details of \texorpdfstring{\ours{}}{Arbor}}
\label{app:arbor-details}

Figure~\ref{fig:appendix-overall-framework} summarizes the implementation-level
agent internals used by \ours{}. 

\begin{figure}[!t]
  \centering
  \includegraphics[width=\linewidth]{fig/appendix_workflow.pdf}
  \caption{Agent-level implementation details of \ours{}.}
  \label{fig:appendix-overall-framework}
\end{figure}

\subsection{Prompts}
\label{app:prompts}

\subsubsection{Coordinator Prompt}
\label{app:meta-agent-prompt}

\begin{promptbox}[title=\textbf{Coordinator Prompt}]
\noindent\textbf{Role.}~You are the persistent coordinator in \arbor{}.
Your job is to maintain the hypothesis tree as the shared research state,
choose which hypotheses deserve execution, and convert experimental feedback
into reusable evidence. You do not edit the target artifact directly. All
implementation work is delegated to short-lived Executors, while you own the
tree, the research frontier, and the merge/prune/stop decisions.

\smallskip
\noindent\textbf{Runtime and durable state.}~You run in a single persistent
ReAct loop. When the conversation approaches the context limit, older turns may
be compressed; the on-disk hypothesis tree is therefore the source of truth.
Every controlled mutation is saved as JSON for tools and regenerated as
Markdown for human inspection. Operator messages prefixed with
\texttt{[user note]} are instructions or questions that must be handled at the
next safe point before launching new Executors or merging branches.

\smallskip
\noindent\textbf{Research contract.}~At initialization, inspect the target
codebase and record the AO contract via \texttt{TreeSetMeta}: the objective,
metric direction, development evaluator $\Eval_\dev$, held-out evaluator
$\Eval_\test$, dataset paths, baseline score, and evaluation commands. Commands
must use the template variables \texttt{\{cwd\}} and \texttt{\{node\_id\}} so
that the same contract can be injected safely into isolated executor worktrees.

\smallskip
\noindent\textbf{Hypothesis tree.}~The root node represents the initial
artifact. Depth-1 children are broad research directions; deeper nodes are more
specific refinements, alternatives, or corrections under their parent. Each
node stores a lifecycle status (\texttt{pending}, \texttt{running},
\texttt{done}, \texttt{merged}, or \texttt{pruned}), a dev score, a factual
result, a structured insight, and a git branch reference \texttt{code\_ref}.
The tree is not a transcript of tool calls: it is the compact record of what
was tried, what happened, why it happened, and how that evidence should shape
future hypotheses.

\smallskip
\noindent\textbf{Coordinator loop.}~Repeat the following cycle until the
budget is exhausted, no promising frontier remains, or the objective has been
satisfied.
\begin{enumerate}[label=\textbf{\arabic*.},leftmargin=1.6em,itemsep=2pt,topsep=2pt]
  \item \textbf{Observe.} Read the current tree with \texttt{TreeView}; inspect
  the current best artifact, recent executor reports, and experiment logs when
  needed. Reconstruct the live research frontier from the persisted tree rather
  than from memory alone.
  \item \textbf{Constrain.} Use \texttt{TreeView(format=``constraints'')} to
  load tree shape, root insight, pruned lessons, and validated findings. Treat
  pruned lessons as negative constraints and validated findings as assumptions
  to build on.
  \item \textbf{Ideate.} Select a parent node and propose a small set of
  executable child hypotheses. Each candidate added with \texttt{TreeAddNode}
  must state its mechanism, hypothesis, observable success condition, and known
  conflicts or risks. Do not add candidates that merely repeat a pruned failure
  mode.
  \item \textbf{Select.} Choose pending leaves for the next batch using
  accumulated evidence, expected impact, implementation cost, and diversity
  across active directions. Prefer hypotheses that can teach the tree something
  useful even when they fail.
  \item \textbf{Execute.} Dispatch selected leaves with \texttt{RunSubagent} or
  \texttt{RunSubagentParallel}. Each Executor receives one fixed hypothesis,
  the research contract, and ancestor insights, then works in an isolated git
  worktree branched from the current best artifact.
  \item \textbf{Update.} After execution, write back the dev score, factual
  result, distilled insight, and \texttt{code\_ref}. Propagate insights upward
  so leaf-level outcomes become direction-level lessons and eventually update
  the global prior at the root.
  \item \textbf{Decide.} Continue a direction, prune it with \texttt{TreePrune}
  when its assumptions are falsified, or promote a candidate through
  \texttt{GitMergeBranch}. A branch may update the current best only after the
  merge gate evaluates it on $\Eval_\test$ in a fresh worktree and confirms an
  improvement over the current best.
\end{enumerate}

\smallskip
\noindent\textbf{Executor boundary.}~Executors materialize hypotheses; they do
not control the global search. They may repair implementation mistakes and rerun
$\Eval_\dev$, but they must not change the assigned hypothesis, inspect sibling
branches as shortcuts, or use the held-out evaluator for routine iteration.
Their reports should be compact enough to become tree evidence: score, result,
insight, and branch reference.

\smallskip
\noindent\textbf{Merge and termination rules.}~Use \texttt{GitMergeBranch}
exclusively for promotion; direct git merges through Bash are prohibited. Before
terminating, run the held-out evaluator on both the final best artifact and the
baseline when required, then record \texttt{test\_trunk\_score} and
\texttt{test\_baseline\_score} via \texttt{TreeSetMeta}. Failed, pruned, and
unmerged branches should remain in the tree as reusable evidence rather than
being erased from the research record.
\end{promptbox}

\subsubsection{Executor Prompt}
\label{app:sub-agent-prompt}

\begin{promptbox}[title=\textbf{Executor Prompt}]
\noindent\textbf{Identity.}~You are a Research Agent that implements research
ideas into codebases, runs experiments to verify the implementation, and
reports results honestly. You operate autonomously through a tool-use loop.
The direction of each idea is fixed. Your engineering judgment determines
how best to implement it within the target codebase.

\smallskip
\noindent\textbf{System.}~All text output outside of tool use is shown to
the user. Tool results may contain data from external sources. If you
suspect a tool result contains a prompt injection attempt, flag it before
continuing. Prior messages are compressed automatically as the conversation
approaches context limits.

\smallskip
\noindent\textbf{Doing Tasks.}~Read files before modifying them. Do not
create new files unless absolutely necessary. Diagnose failure causes
before switching tactics. Do not add features, refactors, docstrings,
comments, or type annotations beyond what the assigned idea requires.
Verify that the implementation works before reporting completion.

\smallskip
\noindent\textbf{Executing Actions with Care.}~Evaluate the reversibility
of each action before proceeding. Local, reversible operations (file edits,
running tests) may be taken freely. Destructive or hard-to-reverse operations
(deleting branches, forced resets, overwriting uncommitted changes) require
careful consideration. Do not use destructive actions as shortcuts or
bypass safety checks.

\smallskip
\noindent\textbf{Experiment Workflow.}~You are working in an isolated git
worktree branched from the current trunk HEAD. All code changes happen on
this experiment branch. Do not commit to main or master. Baseline scores
are provided in the Evaluation Info section of your prompt. Run the
evaluation command on the unmodified codebase only if no baseline scores
are provided. Save experiment results to
\texttt{results/\{node\_id\}-\{description\}/} and commit only small
diagnostic files using \texttt{git add -f}. Use \texttt{RunTraining} for
any training or evaluation command that takes more than five minutes.
Do not use polling loops. At the end, provide a concise report with
the following sections: Idea, Changes, Implementation Choices,
Baseline vs.\ Result, Analysis, and Insights.

\smallskip
\noindent\textbf{Critical Rules.}~The idea direction is non-negotiable.
You may not substitute a fundamentally different approach. Implementation
choices (specific architecture, hyperparameters, code placement) are yours.
Report significant implementation choices explicitly. If the idea
underperforms after a good-faith implementation, that is the finding to report.
\end{promptbox}

\subsection{Algorithm Workflow}
\label{app:algorithm-flow}

Algorithm~\ref{alg:appendix-htr} expands the HTR pseudocode from the main
paper (Algorithm~\ref{alg:htr}) with implementation-level detail. Notation
follows the main paper: $\iota_n$ denotes the insight recorded at node $n$,
$b_n$ the git branch reference, $s_n$ the $\Eval_\dev$ score, $r_n$ the
factual result, and $\iota_{\mathrm{anc}(n)}$ the concatenated insights on
the path from $n_0$ to $n$'s parent.

\begin{algorithm}[!t]
\small
\DontPrintSemicolon
\SetKwInOut{Input}{Input}\SetKwInOut{Output}{Output}
\SetKwProg{Fn}{Procedure}{:}{}
\SetKwFunction{Exec}{Executor}
\caption{Hypothesis Tree Refinement (HTR) with expanded implementation detail.
  The coordinator owns $\Tree$, and each \textsc{Executor} owns one worktree.}
\label{alg:appendix-htr}
\Input{$\mathcal{P}=(\Mat_0,\Obj,\Eval_{\dev},\Eval_{\test})$, budget $B$, branching $k$, parallelism $P$}
\Output{best artifact $\Mat^\star$, annotated hypothesis tree $\Tree$, and run summary report}

\tcp{Initialization}
init $\Tree=(\{n_0\},\emptyset)$, $b_{n_0}\leftarrow\Mat_0$, $\Mat_{\mathrm{best}}\leftarrow\Mat_0$\;
run $\Eval_{\dev}(\Mat_0)$ and record baseline score and eval command in $\Tree.\mathrm{meta}$ via \texttt{TreeSetMeta}\;

\tcp{Main coordinator loop}
\While{$B$ left $\wedge$ pending leaves exist $\wedge$ no stop signal}{
  \tcp{Step 1: OBSERVE, build constraint view}
  $\mathcal{V}\leftarrow\textsc{Observe}(\Tree,\Mat_{\mathrm{best}})$
    \tcp*{shape, root insight, pruned lessons, validated findings}

  \tcp{Step 2: IDEATE, skill-gated hypothesis proposal}
  $\mathcal{V}_c \leftarrow \texttt{TreeView}(\mathrm{format}{=}\texttt{constraints})$
    \tcp*{hard constraints: no re-tread of pruned directions}
  load \texttt{idea\_drafting} skill via \texttt{LoadSkill}
    \tcp*{mandatory gate: must precede any candidate proposal}
  $p\leftarrow\textsc{Select}_{\textsf{parent}}(\mathcal{V})$\;
  \ForEach{surviving candidate after Fatal-Flaw Scan}{
    attach pending child $n^{(i)}$ with hypothesis $h^{(i)}$ = (Mechanism, Hypothesis, Observable, Conflicts)\;
    $\iota_{\mathrm{anc}(n^{(i)})}\leftarrow$ insights on $\textsf{path}(n_0\to p)$
      \tcp*{injected into Executor prompt}
    prune \texttt{idea\_drafting} scratch from coordinator context
      \tcp*{elide skill body + reasoning post-TreeAddNode}
  }

  \tcp{Step 3: SELECT frontier for parallel dispatch}
  $L\leftarrow$ up to $P$ pending leaves under $\textsc{Select}_{\textsf{frontier}}(\mathcal{V})$\;

  \tcp{Step 4: DISPATCH, parallel Executor dispatch}
  $\{(s_n,r_n,\iota_n,b_n)\}_{n\in L}\leftarrow\textbf{parallel }\Exec(h_n,\iota_{\mathrm{anc}(n)},\Mat_{\mathrm{best}})$\;

  \tcp{Step 5: UPDATE, write back and propagate}
  \ForEach{$n\in L$, $a\in\textsf{path}(n_0\to n)$}{
    write back $(s_n,r_n,\iota_n,b_n)$ to node $n$ in $\Tree$ and set $n.\mathrm{status}\leftarrow\texttt{done}$\;
    $\iota_a\leftarrow\textsc{Abstract}(\{\iota_c\}_{c\in\textsf{ch}(a)})$
      \tcp*{propagate insights upward to root}
  }
  check convergence: inject intervention if $\geq w$ consecutive non-improving experiments\;

  \tcp{Step 6: DECIDE, merge gate or prune}
  $n^\dagger\leftarrow\arg\max_{n\in L}\,s_n$\;
  \If{$s_{n^\dagger}$ exceeds best score by $\geq\theta$}{
    create detached worktree at $b_{n^\dagger}$ and run $\Eval_{\test}$ with template substitution\;
    \lIf{$S_{\test}(b_{n^\dagger})>S_{\test}(\Mat_{\mathrm{best}})$}{
      $\Mat_{\mathrm{best}}\leftarrow\textsf{merge}(b_{n^\dagger})$ and update $\Tree.\mathrm{meta}.\mathrm{trunk\_score}$
    }
  }
  prune subtrees falsified by $\{\iota_n\}_{n\in L}$ and persist $\Tree$ to JSON + Markdown\;
}
run $\Eval_{\test}(\Mat_{\mathrm{best}})$ and record \texttt{test\_trunk\_score} and \texttt{test\_baseline\_score}\;
\Return{$\Mat^\star\leftarrow\Mat_{\mathrm{best}}$, $\Tree$, run summary}\;

\BlankLine
\Fn{\Exec{$h_n,\,\iota_{\mathrm{anc}(n)},\,\Mat_{\mathrm{best}}$}}{
  branch name $\leftarrow \textsf{slug}(\mathrm{node\_id}) + \textsf{slug}(h_n) + \mathrm{SHA1}(h_n)_{[:8]}$\;
  create worktree $W_n$ in $\texttt{/tmp/}$ from current best branch HEAD on a fresh branch\;
  inject eval command (with \texttt{\{cwd\}$\to W_n$}, \texttt{\{node\_id\}$\to n$}) and $\iota_{\mathrm{anc}(n)}$ into prompt\;
  \Repeat{run ok $\wedge$ $h_n$-path exercised, or turn cap reached}{
    $\Delta\leftarrow\textsf{Implement}(h_n,\iota_{\mathrm{anc}(n)},W_n)$\;
    $(s_n,r_n)\leftarrow\Eval_{\dev}(\textsf{apply}(\Delta,W_n))$
      \tcp*{repair $\Delta$ only, direction $h_n$ is fixed}
  }
  filter $\Delta$: commit implementation files only, skip logs/checkpoints/caches\;
  remove worktree directory and retain branch $b_n$ for later merge gate\;
  \Return{$(s_n,\;r_n,\;\textsc{Distill}(h_n,\Delta,r_n),\;b_n)$}\;
}
\end{algorithm}

\subsection{Key Design of \arbor{} Framework}
\label{app:framework-key-design}
General-purpose coding agents such as Codex and Claude Code are built for
general-purpose software tasks: they chain tool calls on a single working tree
to edit, test, and fix code against a goal that is already well specified.
\arbor{} instead targets the auto-research setting, and we make a series of
engineering choices that specialize the agent for it so that the system can
flexibly adapt to different research needs and manage many experiments over a
long horizon. We group its key engineering designs into four areas: (i)~hypothesis-tree management,
(ii)~experiment management, (iii)~long-horizon operation, and (iv)~functional
extensibility through skills and plugins.

\paragraph{Hypothesis-tree management.}
Unlike a code agent whose memory is the linear chat transcript, \arbor{}
externalizes its research state into an explicit hypothesis tree (the
\emph{idea tree}), an in-memory object that is the single authoritative record
of a run. Each node holds a hierarchical dotted address (e.g.\ \texttt{ROOT},
\texttt{1}, \texttt{1.1}) that encodes its path from the root, its
\texttt{parent\_id} and \texttt{children\_ids}, a once-written
\texttt{hypothesis}, a \texttt{status} field tracing the lifecycle
\texttt{pending}\,$\to$\,\texttt{running}\,$\to$\,\texttt{done}\,$\to$\,%
\{\texttt{merged},\,\texttt{pruned}\}, a development \texttt{score}, a factual
\texttt{result}, a distilled \texttt{insight}, and a \texttt{code\_ref} branch
pointer to the artifact. Depth-1 nodes are broad directions and deeper nodes are
concrete refinements, so edges encode hypothesis refinement rather than
chronological actions. The coordinator never touches the raw structure; it
operates the tree only through a small set of typed tools
(\texttt{TreeAddNode}, \texttt{TreeUpdateNode}, \texttt{TreePrune},
\texttt{TreeSetMeta}, \texttt{TreePropagate}) and read projections
(\texttt{TreeView}). Storing an experiment is thus a controlled mutation: the
executor's structured report is parsed and written into the node's
\texttt{score}/\texttt{result}/\texttt{insight}/\texttt{code\_ref} fields, and
the tree is serialized to JSON (and rendered to Markdown) after \emph{every}
mutation. Crucially, when a node finishes \arbor{} \emph{backpropagates} its
insight: \texttt{propagate\_insights} walks from the node's parent to the root
and, at each ancestor, an LLM synthesizes the insights of that ancestor's
children into a concise ($<$200-word) summary. Leaf insights describe concrete
implementations, parent insights summarize families of interventions, and the
root insight maintains a global understanding of the problem. This layerwise
abstraction lets the coordinator reason at the right granularity without
rereading raw logs. The concrete schema and persistence format are detailed in
Section~\ref{app:idea-tree-storage}.

\paragraph{Experiment management.}
Each pending node is executed by an executor dispatched through
\texttt{RunSubagent} or \texttt{RunSubagentParallel}. Rather than editing the
shared working tree in place, \arbor{} creates a fresh git worktree branched
from the current trunk \texttt{HEAD} for each candidate, so every hypothesis gets
a clean, independently recoverable experimental boundary and several executors
can edit overlapping files concurrently without corrupting the trunk or one
another. At launch the executor is injected with the assigned hypothesis, the
ancestor insights along its path, and the task
objective, development evaluator $\Eval_\dev$, held-out evaluator $\Eval_\test$,
metric direction, data split, and baseline score stored in the tree metadata
via \texttt{TreeSetMeta}. Its permission boundary is deliberately narrow: the
executor's tools (Table~\ref{tab:sub-agent-tools}) act \emph{only} within its
own worktree. It cannot read the trunk, inspect sibling branches, mutate the
tree, or run $\Eval_\test$. It may repair implementation bugs and choose
reasonable engineering details, but it may not swap the assigned hypothesis for
a different research direction, which keeps the returned evidence attributable
to the node actually tested. The executor reports back four separated fields:
the development score $s_n$, a factual \texttt{result} $r_n$ (raw observations
such as errors, curves, and metric breakdowns), a distilled \texttt{insight}
$\iota_n$ (the causal lesson), and the branch reference $b_n$. The coordinator
auto-extracts these, updates the node, and decides admission: a result counts as
a \emph{useful gain} only when it improves $\Eval_\dev$ over the trunk by at
least the merge threshold, which triggers the \texttt{GitMergeBranch} held-out
gate that re-runs $\Eval_\test$ in a separate detached worktree and promotes the
artifact only if it strictly beats the current best. Every other outcome is
treated as \emph{failure evidence} rather than noise: the node's insight and, if
the direction is abandoned, its \texttt{TreePrune} reason are aggregated into
the constraint view (\texttt{TreeView(format="constraints")}) that conditions
the next \textsc{Ideate} step, so negative results actively narrow the search
space. Tool-level details of dispatch, score extraction, template-variable
substitution (\texttt{\{cwd\}}, \texttt{\{node\_id\}}), and artifact filtering
are given in Section~\ref{app:agent-tools-params}.

\paragraph{Long-horizon operation.}
A single auto-research run can span hundreds of turns and many hours of
evaluation, far beyond one context window, so \arbor{} adds explicit mechanisms
to keep the search productive instead of stalling on early failures or chasing
noisy evaluation swings. First, because all durable state lives in the persisted
idea tree rather than the transcript, the run survives crashes, agent restarts,
and context compression; post-commit context pruning further elides spent
\textsc{Ideate} scratch work and loaded skill bodies once a candidate is
committed, bounding context growth. Second, long training or evaluation commands
are routed through \texttt{RunTraining}, which blocks until completion or timeout
while continuously capturing partial metrics, progress logs, and checkpoints, so
even a run that times out at 80\% of its epochs returns actionable evidence
instead of a silent failure. Third, a \emph{convergence detector} monitors
recent score velocity over a sliding window and counts consecutive
non-improving experiments, escalating through \texttt{warn}, \texttt{paradigm\_shift},
and \texttt{stop} signals; it also flags \emph{parent exhaustion} when a parent's
recent children all fail to beat the trunk, prompting the coordinator to
summarize the failure pattern and open a fresh depth-1 direction rather than
over-exploring a dead branch. A meaningful-improvement threshold prevents a
single noisy uptick from resetting this signal, balancing premature stagnation
against unbounded exploration.

\paragraph{Functional extensibility (skills and plugins).}
To remain adaptable across research domains without changing the core loop,
\arbor{} exposes two extension surfaces. \emph{Skills} are markdown documents
with YAML frontmatter, discovered by a registry from both the built-in directory
and a project-local \texttt{.research\_agent/skills/} override, and loaded on
demand via \texttt{LoadSkill}. The \textsc{Ideate} protocol, for instance, loads
\texttt{idea\_drafting}, \texttt{first\_principles\_probe}, and
\texttt{fatal\_flaw\_scan} before proposing candidates, then prunes their bodies
from context after each commit, so reasoning guidance is injected just-in-time
and does not permanently inflate context. \emph{Plugins} are YAML domain
adapters that specialize the system declaratively rather than through code: each
plugin can inject domain guidance at six prompt points (coordinator/executor
init, ideate, decide, preamble, and workflow), declare an evaluation contract,
mark protected paths and required outputs for the merge guard, override runtime
configuration through named profiles, and tune convergence thresholds. For
example, the \texttt{mle\_kaggle} plugin configures a Kaggle/MLE-bench task with
its metric direction, evaluation command, protected data paths, and time-budget
profiles entirely in YAML. Together, skills and plugins let \arbor{} adapt to new
artifacts and evaluation regimes while keeping the hypothesis-tree machinery and
agent contracts unchanged.

\subsection{Agent Tools and Hyperparameter Settings}
\label{app:agent-tools-params}

\subsubsection{Coordinator Tools}
\label{app:meta-agent-tools}

The coordinator's tool set reflects a strict division of labor: it may read
any file in the repository and inspect the tree in any projection, but it
may never edit code directly. All implementation work is delegated to
executors via \texttt{RunSubagent} or \texttt{RunSubagentParallel}. This
design is intentional. If the coordinator could edit code directly, it
would be tempted to make small local fixes without creating a new tree
node, which would break the invariant that every code change is traceable
to a hypothesis. Enforcing the edit boundary at the tool level makes this
invariant a system property rather than a behavioral expectation.

Table~\ref{tab:meta-agent-tools} lists the full tool set. The tree tools
(\texttt{TreeView}, \texttt{TreeAddNode}, \texttt{TreeUpdateNode},
\texttt{TreePrune}, \texttt{TreeSetMeta}, \texttt{TreePropagate}) are the
primary interface through which the coordinator manages the shared research
state. The dispatch tools (\texttt{RunSubagent},
\texttt{RunSubagentParallel}) are the boundary through which the coordinator
hands off implementation to executors and receives structured evidence back.
The merge tool (\texttt{GitMergeBranch}) is the only path through which a
candidate branch can be promoted to the current best, enforcing the held-out gate at the
tool level rather than relying on the coordinator to remember to verify
before merging. Table~\ref{tab:meta-agent-hyperparams} reports the key
hyperparameters used across all experiments.

\begingroup
\footnotesize
\setlength{\tabcolsep}{3pt}
\renewcommand{\arraystretch}{1.02}
\setlength{\LTpre}{4pt}
\setlength{\LTpost}{6pt}
\begin{xltabular}{\linewidth}{>{\raggedright\arraybackslash}p{3.8cm}>{\raggedright\arraybackslash}X}
\caption{Tools available to the coordinator.}
\label{tab:meta-agent-tools}\\
\toprule
\textbf{Tool} & \textbf{Description} \\
\midrule
\endfirsthead
\caption[]{Tools available to the coordinator (continued).}\\
\toprule
\textbf{Tool} & \textbf{Description} \\
\midrule
\endhead
\midrule
\multicolumn{2}{r}{\footnotesize continued on next page}\\
\endfoot
\bottomrule
\endlastfoot
\texttt{TreeView} & Inspect the idea tree in five formats: \texttt{compact} (status overview), \texttt{full} (Markdown rendering), \texttt{node} (single-node detail), \texttt{pending} (pending leaf list), and \texttt{constraints} (aggregated pruned lessons and validated findings used as hard constraints in the IDEATE phase). \\
\texttt{TreeAddNode} & Add a child node to the tree with a four-field hypothesis block (Mechanism, Hypothesis, Observable, Conflicts). Triggers IDEATE context pruning after each commit. \\
\texttt{TreeUpdateNode} & Update mutable fields of an existing node (\texttt{status}, \texttt{insight}, \texttt{score}, \texttt{result}, \texttt{code\_ref}, \texttt{related\_work}). \\
\texttt{TreePrune} & Mark a node and its subtree as pruned with a written reason. The reason is aggregated into the constraints block for future IDEATE rounds. \\
\texttt{TreeSetMeta} & Write evaluation metadata to the tree root: evaluation commands (\texttt{eval\_cmd}, \texttt{eval\_cmd\_test}), dataset paths, baseline score, best score, and test scores. Automatically injected into every Executor's prompt. \\
\texttt{TreePropagate} & Re-propagate insights from a node upward to the root after manual corrections via \texttt{TreeUpdateNode}. \\
\texttt{RunSubagent} & Dispatch a single Executor to implement and evaluate a pending tree node. Creates an isolated git worktree from trunk HEAD, injects ancestor insights and evaluation metadata, auto-extracts results from the structured report, and updates the tree node. \\
\texttt{RunSubagentParallel} & Dispatch two to four Executors concurrently on independent tree nodes. Each Executor runs in its own isolated worktree. \\
\texttt{GitMergeBranch} & Validate and promote a candidate branch to the current best. Creates a temporary detached worktree at the source branch, runs $\Eval_\test$ with score extraction, and merges only if the validated test score exceeds the current best score. Protected branches (main, master) cannot be merge targets. \\
\texttt{Bash} & Execute shell commands for read-only codebase inspection and environment queries. The coordinator never uses Bash to edit code. \\
\texttt{FileRead} & Read the contents of a file within the target repository. \\
\texttt{Grep} & Search for a text pattern across repository files. \\
\texttt{Glob} & Find files matching a glob pattern within the repository. \\
\texttt{LoadSkill} & Load a skill document on demand from the registry. In the IDEATE protocol, \texttt{idea\_drafting}, \texttt{first\_principles\_probe}, and \texttt{fatal\_flaw\_scan} are loaded before proposing candidates. Loaded content is pruned from context after each \texttt{TreeAddNode} commit. \\
\texttt{SearchIdeaContext} & Dispatch a SearchAgent in the background to annotate a tree node with related work. Returns immediately. The SearchAgent runs concurrently with ongoing IDEATE and dispatch work and writes a structured annotation (summary, novelty assessment, related papers) to \texttt{node.related\_work} when finished. \\
\texttt{SearchIdeaContextParallel} & Dispatch multiple background SearchAgents concurrently for a list of node IDs. \\
\texttt{SearchStatus} & Return the count of in-flight background SearchAgent tasks. \\
\end{xltabular}
\endgroup

\begin{table}[!htbp]
\centering
\small
\setlength{\tabcolsep}{4pt}
\renewcommand{\arraystretch}{1.08}
\caption{Key hyperparameters for the \arbor{} coordinator and executors.}
\label{tab:meta-agent-hyperparams}
\begin{tabularx}{\linewidth}{>{\raggedright\arraybackslash}p{4.8cm}>{\raggedright\arraybackslash}p{2.4cm}>{\raggedright\arraybackslash}X}
\toprule
\textbf{Parameter} & \textbf{Default} & \textbf{Description} \\
\midrule
Backbone model & Claude Opus 4.6 & LLM for both coordinator and executors (configurable per role) \\
Max cycles & 20 & Hard cap on total completed experiments \\
Max coordinator turns & 500 & Maximum ReAct turns for the persistent coordinator loop \\
Executor max turns & 50 & Maximum turns per executor \\
Executor timeout & 48\,h & Wall-clock limit per executor \\
Context window & 200\,K tokens & Coordinator context limit \\
Compression threshold & 0.80 & Compress context when it reaches 80\% of the context window \\
Kept recent turns & 20 & Number of most recent turns preserved after each compression \\
Merge threshold & 5.0\% & Minimum $\Eval_\dev$ improvement required to invoke the merge gate \\
$\Eval_\test$ evaluation retries & 1 & Additional held-out evaluation attempts after a transient failure \\
Convergence window & 5 & Sliding window size for score velocity computation \\
Convergence warn / stop & 3 / 8 & Consecutive non-improving experiments to trigger warning / stop \\
\bottomrule
\end{tabularx}
\end{table}

\subsubsection{Executor Tools}
\label{app:sub-agent-tools}

The executor's tool set is designed around a single constraint: every action
must be local to its worktree. The executor has no tools for reading the
current best, inspecting sibling branches, or modifying the shared tree. This
isolation is not merely a safety measure. It is what allows multiple
executors to run concurrently on overlapping codebases without coordination
overhead. Each executor operates as if it has an exclusive copy of the
repository, because its worktree is literally a separate checked-out
directory that shares the git object store but maintains its own HEAD and
index. The \texttt{RunTraining} tool deserves special mention: for experiments
that involve multi-hour training runs, using \texttt{Bash} with a timeout
would either cut the run short or block the executor for an unbounded
duration. \texttt{RunTraining} instead blocks until the command terminates
or the configured timeout is reached, continuously capturing intermediate
metrics, progress logs, and checkpoints so that even a timed-out run
returns actionable partial evidence. This matters in practice: a training
run that completes 80\% of its epochs before timing out often reveals enough
about a direction's trajectory to inform a principled prune or continuation
decision, which is more useful than a silent timeout that leaves the
coordinator without evidence.

\begin{table}[!htbp]
\centering
\small
\setlength{\tabcolsep}{4pt}
\renewcommand{\arraystretch}{1.08}
\caption{Tools available to each executor.}
\label{tab:sub-agent-tools}
\begin{tabularx}{\linewidth}{>{\raggedright\arraybackslash}p{2.8cm}>{\raggedright\arraybackslash}X}
\toprule
\textbf{Tool} & \textbf{Description} \\
\midrule
\texttt{Bash} & Execute shell commands in the worktree with configurable per-call timeouts (default 600\,s, maximum 86{,}400\,s). \\
\texttt{RunTraining} & Execute a long-running training or evaluation command and block until completion or timeout. Automatically captures partial metrics, progress logs, and checkpoints. Supports staged budgets (smoke, pilot, full) with configurable wall-times. Maximum timeout is 604{,}800\,s (seven days). \\
\texttt{FileRead} & Read the full contents of a file within the worktree. \\
\texttt{FileEdit} & Edit a file by replacing an exact string match with new content, enabling surgical edits that minimize unintended side effects. \\
\texttt{FileWrite} & Write the complete contents of a file, creating or overwriting it. \\
\texttt{Grep} & Search for a text or regex pattern across files in the worktree. \\
\texttt{Glob} & Find files matching a glob pattern within the worktree. \\
\texttt{SubAgent} & Spawn a nested executor with a custom prompt for modular decomposition of complex implementation tasks. The nested agent inherits the same tool set and worktree context. \\
\bottomrule
\end{tabularx}
\end{table}

\subsubsection{Evaluation and Merge Tools}
\label{app:evaluation-merge-tools}

A recurring challenge in multi-agent systems is ensuring that evaluation
commands measure the right code. In \arbor{}, each executor runs in a
different directory, on a different branch, with potentially different data
locations. A naively hardcoded evaluation command would silently evaluate
the wrong artifact. The template variable system addresses this by requiring
the coordinator to specify evaluation commands using two placeholders:
\texttt{\{cwd\}} is substituted with the executor's worktree path before
the command reaches the executor's prompt, and \texttt{\{node\_id\}} is
substituted with the tree node identifier so that results from concurrent
experiments are written to distinct output locations and never overwrite
one another. The coordinator is also explicitly prohibited from hardcoding
absolute paths in evaluation commands. Only \texttt{\{cwd\}}-relative
paths are permitted.

Score extraction from evaluator output follows a two-stage pipeline
that handles the diversity of real evaluation scripts. Many harnesses
already emit a structured JSON block with a \texttt{score} key. The system
attempts to parse this first. If no valid JSON block is found, as is
common with training scripts that report loss or accuracy in free-form
log lines, a secondary LLM call reads the full command output and
extracts the primary metric as a percentage. The distinction between a
primary metric and ancillary logging is task-specific, so the LLM call
uses a concise system prompt that instructs it to return only a single
JSON object and to pick the most prominent performance figure if multiple
metrics appear.

The held-out evaluator $\Eval_\test$ is architecturally separated from
the development loop. Executors are instructed never to run $\Eval_\test$.
It is accessible only through the \texttt{GitMergeBranch} tool, which
creates a fresh detached worktree at the candidate branch's HEAD, applies
the same template substitution, and runs \texttt{eval\_cmd\_test} in
complete isolation from the main repository. Configurable retry logic
with exponential backoff handles transient infrastructure failures such as
gpu allocation timeouts or flaky network reads. The merge is admitted only
if the extracted test score strictly exceeds the current best score under
the metric direction specified in the task. This two-worktree design uses one
worktree for development evaluation inside the executor and one for held-out
evaluation inside the merge gate. It ensures that neither the executor nor
the coordinator
ever evaluates the candidate artifact in the same directory as the current best,
eliminating any risk of result leakage between branches.

Before a worktree branch is committed, \arbor{} applies artifact filtering
to separate implementation changes from generated outputs. The system
inspects the full worktree diff and classifies changed paths as either
implementation files or artifacts. Raw logs, model checkpoints,
cache directories, generated data, and large binary files are excluded
from the commit, keeping experiment branches compact and ensuring that
subsequent three-way merges into the current best branch involve only meaningful code
differences. This filtering is conservative by design: a borderline file
is included rather than excluded, because a missing implementation file
breaks the merge, whereas a spurious large file only adds storage overhead.

\subsection{Idea Tree Data Structure and Storage}
\label{app:idea-tree-storage}

The implementation stores \arbor{}'s hypothesis tree as an idea-tree object.
This object is the primary durable data structure of a run: it records the
research frontier, completed experiments, rejected directions, accepted
artifacts, and reusable insights in one inspectable state.

\paragraph{Logical schema.}
The root node anchors the initial artifact and task contract. Depth-1 nodes
represent broad research directions, while deeper nodes represent concrete
refinements, alternatives, or corrections under their parent. Edges therefore
encode hypothesis refinement rather than chronological agent actions. During a
run, pending leaves form the executable frontier; completed, merged, and pruned
nodes remain in the tree as evidence for future proposal and selection.
Each node carries the fields listed in Table~\ref{tab:node-fields}.

\begin{table}[ht]
\centering
\small
\setlength{\tabcolsep}{4pt}
\renewcommand{\arraystretch}{1.10}
\caption{Fields of a single node in the idea tree.
The first four fields form the structural skeleton;
the remaining fields are populated progressively as the node
moves through its status lifecycle.}
\label{tab:node-fields}
\begin{tabularx}{\linewidth}{>{\raggedright\arraybackslash}p{2.4cm}%
>{\raggedright\arraybackslash}p{2.8cm}>{\raggedright\arraybackslash}X}
\toprule
Field & Type & Description \\
\midrule
\texttt{id} & string
  & Hierarchical address encoding position in the tree
    (e.g.\ \texttt{"ROOT"}, \texttt{"1"}, \texttt{"1.1"}).
    The dot-separated prefix gives the unique path from root. \\
\texttt{parent\_id} & string\,\textbar\,null
  & Identifier of the parent node. \texttt{null} for the root. \\
\texttt{children\_ids} & list[string]
  & Ordered identifiers of child nodes, appended as the
    coordinator expands the tree during \textsc{Ideate}. \\
\texttt{depth} & int
  & Integer depth: $0$~=~root, $1$~=~direction node,
    $2$~=~branch node.
    The maximum depth is governed by \texttt{max\_depth}. \\
\addlinespace
\texttt{hypothesis} & string
  & Research hypothesis assigned before execution.
    Written once by the coordinator during \textsc{Ideate} and
    never modified thereafter. \\
\texttt{status} & \{pending, running,\newline done, merged, pruned\}
  & Status lifecycle: \texttt{pending}~$\to$~\texttt{running}~$\to$~%
    \texttt{done}, then \texttt{done}~$\to$~\texttt{merged}
    if the merge gate passes, or \texttt{done}~$\to$~\texttt{pruned}
    otherwise. \\
\texttt{score} & float\,\textbar\,null
  & Absolute scalar score on $\Eval_{\dev}$ as reported by the executor.
    Null until the node has been executed and scored. \\
\texttt{result} & string
  & Factual record of the experiment outcome, written by the executor
    as a direct observational report before any interpretation is applied. \\
\texttt{insight} & string
  & Structured lesson extracted by \textsc{Distill} from the outcome,
    then backpropagated and aggregated at each ancestor up to the root. \\
\texttt{code\_ref} & string\,\textbar\,null
  & Git branch name pointing to the implementation artifact produced
    by the executor.
    Null if no artifact was committed. \\
\bottomrule
\end{tabularx}
\end{table}

\paragraph{Persistent storage.}
The idea tree is serialized to JSON after every controlled mutation and serves
as the sole authoritative record of research state accumulated during a run.
No hypothesis, score, artifact reference, or lesson exists outside this
structure. The same state is also rendered to Markdown for dashboards and human
inspection. The JSON representation is consumed by tools, while the Markdown
view provides a readable account of the research process. Since every mutation
is persisted immediately, the run can recover from crashes, agent restarts, and
context-window compression without relying on conversational memory.

\paragraph{Agent-facing access.}
Agents do not edit the raw JSON file directly. The coordinator accesses the
tree through controlled read views and mutation tools: compact frontier views,
single-node views, pending-leaf views, and the constraint view used before
ideation; plus mutations for adding nodes, updating lifecycle fields, pruning
subtrees, setting metadata, and propagating insights. Executors receive only the
assigned hypothesis, the research contract, and ancestor insights. This access
pattern keeps the tree as a shared state object rather than a low-level log of
tool traces.

\FloatBarrier

\section{Details of the AO Test Suite}
\label{app:ao-suite-details}
\label{app:benchmark-details}

\subsection{Optimizer Design}
\label{app:benchmark-optimizer-design}

NanoGPT-Bench~\citep{nanogptbench} Track~3 is a collaborative benchmark
for discovering efficient neural network optimizers.
Unlike the main NanoGPT speedrun, which minimizes wall-clock time by any
means, Track~3 targets \emph{step count}: methods that are slow in
wall-clock terms but reach the target loss in fewer steps are valid and
desirable.
All optimizer designs are evaluated on the same fixed architecture, dataset,
and training script, so the only lever available is the optimization
algorithm and its hyperparameters.

\paragraph{Baseline.}
We initialize from the official tuned Muon optimizer setting provided by the
benchmark authors.
This configuration applies Muon to transformer block weights and AdamW to
embeddings, the output projection, and scalar parameters (biases and gains).
With these tuned hyperparameters, the baseline reaches a FineWeb validation
loss of $\leq 3.28$ in 3{,}325 steps.
Selecting a well-tuned official baseline as the starting material is
intentional: it places the agent in a regime where naive hyperparameter
sweeps are unlikely to yield large gains, thereby testing the agent's ability
to identify and implement genuinely novel optimizer improvements rather than
harvesting easy wins from an undertuned starting point.

\paragraph{Development and test split.}
Agents iterate with the development evaluator throughout the search loop.
After the search completes, the selected optimizer is re-evaluated with two
held-out random seeds, and the reported test score is the average step count
across those runs.

\paragraph{Evaluator.}
The evaluator runs the benchmark's official evaluation script
\texttt{run\_eval.py}.
The script launches the training job, monitors validation loss, and
terminates as soon as \texttt{val\_loss} $\leq 3.28$ is first reached.
The score is the step count at termination, with lower being better.
If the target is never reached, a penalty score exceeding 7{,}000 is
assigned.
All evaluation runs are executed on four NVIDIA A100-80\,GB GPUs using the
official benchmark evaluation code.

\paragraph{Agent instruction.}
The agent is initialized with the following task description.

\begin{promptbox}[title=\textbf{Optimizer Design Task Instruction}]
Improve the training-step efficiency of the NanoGPT optimizer on FineWeb.
The training script \texttt{train\_gpt\_simple.py} trains a 124\,M-parameter
GPT-2 on FineWeb using Muon for transformer weights and AdamW for embeddings,
output projection, and scalar parameters.

\smallskip
\noindent\textbf{Goal.}~Minimize the number of training steps required to
reach \texttt{val\_loss}~$\leq 3.28$.
The current baseline achieves this target in 3{,}325 steps.

\smallskip
\noindent\textbf{Evaluation.}~Run \texttt{python run\_eval.py} from the
project directory.
The script launches training, monitors validation loss, and terminates at the
first step where \texttt{val\_loss}~$\leq 3.28$.
The score is the step count at termination, with lower being better.
If the target is never reached, a penalty score exceeding 7{,}000 is assigned.

\smallskip
\noindent\textbf{Constraints.}~Only \texttt{train\_gpt\_simple.py} may be
modified.
The dataset, batch size, and model architecture must remain unchanged.
Adjusting \texttt{train\_steps} in isolation is not a valid improvement.
Multiple forward passes per step are not permitted.
\end{promptbox}

\subsection{Architecture Design}
\label{app:benchmark-architecture-design}

Architecture Design uses the \texttt{autoresearch} benchmark~\citep{autoresearch},
a compact LLM pretraining task designed for closed-loop research on a real
training codebase.
The agent receives a single-file training implementation and must improve the
final validation loss under a fixed wall-clock training budget.
Unlike Optimizer Design, where the model architecture is fixed and only the
optimizer is modified, this task exposes the full training recipe: model shape,
attention pattern, initialization, optimizer hyperparameters, learning-rate
schedules, batch sizing, and training-loop details may all be changed as long
as the benchmark evaluator and data preparation remain fixed.

\paragraph{Baseline.}
The initial material is the default \texttt{autoresearch} repository.
The main artifact is \texttt{train.py}, a single-GPU decoder-only Transformer
pretraining script derived from \texttt{nanochat}.
Data preparation, tokenization, validation-token selection, and evaluation
helpers live outside the editable artifact, primarily in \texttt{prepare.py}.
The development evaluator uses the fixed validation shard provided by the benchmark.
The default model is a roughly 50M-parameter Transformer, and the baseline
reports the final validation bits-per-byte after the fixed training run.

\paragraph{Development and test split.}
During search, agents optimize only against the development evaluator, which
runs the training script on the benchmark's fixed prepared data and reports the
final validation loss.
The held-out score reported in the main experiments is obtained by rerunning
the selected \texttt{train.py} with two held-out random seeds and averaging the
resulting final losses.
This separates the fast, single-run feedback used for hypothesis exploration
from the seed-averaged score used for final comparison.

\paragraph{Evaluator.}
The evaluator executes
\texttt{uv run train.py} from the benchmark repository.
At the end of training the script prints a structured summary including
\texttt{val\_bpb}, \texttt{peak\_vram\_mb}, \texttt{training\_seconds},
\texttt{num\_steps}, and \texttt{mfu\_percent}.
The primary score is the final \texttt{val\_bpb}, with lower values better.
The training script enforces a 300-second training budget; runs that crash,
time out, fail to fit in memory, or do not emit a parseable final
\texttt{val\_bpb} are treated as failed experiments.
The agent may use secondary diagnostics such as step count and peak memory to
interpret failures, but merge decisions are based on the loss metric.

\paragraph{Agent instruction.}
The agent is initialized with the following task description.

\begin{promptbox}[title=\textbf{Architecture Design Task Instruction}]
Follow \texttt{program.md} as the authoritative task specification for the
\texttt{autoresearch} repository.

\smallskip
\noindent\textbf{Goal.}~Minimize the validation bits-per-byte
(\texttt{val\_bpb}) reported by \texttt{uv run train.py}.
Lower is better.

\smallskip
\noindent\textbf{Evaluation.}~Run \texttt{uv run train.py} from the project
directory and extract the final \texttt{val\_bpb}, together with
\texttt{peak\_vram\_mb}, \texttt{training\_seconds}, \texttt{num\_steps},
and \texttt{mfu\_percent} for diagnostics.
The script has a fixed 5-minute training budget.
If a run crashes, runs out of memory, exceeds the allowed wall-clock budget, or
does not produce a parseable final loss, treat it as a failed experiment unless
there is an obvious implementation bug that can be corrected quickly.

\smallskip
\noindent\textbf{Constraints.}~Only \texttt{train.py} may be modified.
Architecture, hyperparameters, optimizer behavior, batch size, model size,
attention pattern, initialization, schedules, and training-loop details are in
scope.
Do not modify \texttt{prepare.py}, especially \texttt{TIME\_BUDGET},
\texttt{MAX\_SEQ\_LEN}, \texttt{EVAL\_TOKENS}, \texttt{make\_dataloader}, or
\texttt{evaluate\_bpb}.
Do not modify data files, tokenizer artifacts, \texttt{pyproject.toml},
\texttt{uv.lock}, or install new dependencies.
\end{promptbox}

\subsection{Terminal-Bench 2.0}
\label{app:benchmark-terminal-bench}

Terminal-Bench~2.0~\citep{terminalbench} is a benchmark for evaluating
terminal agents on realistic command-line tasks executed inside isolated
Docker containers.
Tasks span 16 categories, including software engineering, security,
scientific computing, data science, games, debugging, and others, across
easy, medium, and hard difficulty levels.
The full task pool contains 89 tasks.

\paragraph{Baseline.}
Prior harness-engineering studies may begin from codebase variants that
differ in interface extensibility, tool set, or prompt structure, creating
comparison confounds.
To ensure a fair and reproducible starting point we use the official
terminal-agent codebase (\texttt{terminus-2}) distributed with the
benchmark verbatim, without any modifications and without introducing
additional interfaces.
The backbone model is GPT-5.5.
The agent operates as a ReAct-style loop that issues tmux keystrokes to
a persistent terminal session.

\paragraph{Development and test split.}
We stratify the 89 tasks by difficulty and sample a 36-task development
set and a 53-task held-out test set.
The stratification ensures that both splits have balanced difficulty
coverage, preventing a scenario where the agent is evaluated on an
atypically easy or hard subset.
Agents iterate exclusively on the development set during the research
loop.
The test set is evaluated only once after the search completes, to
report held-out performance.

\paragraph{Evaluator.}
The evaluator is the official Harbor evaluation harness distributed with
Terminal-Bench~2.0.
Both development and test evaluations use 8 concurrent workers.
The evaluator reports pass rate as the fraction of tasks solved correctly,
with higher being better.
The baseline achieves 58.33\% on the development set (21 out of 36 tasks).

\paragraph{Agent instruction.}
The agent is initialized with the following task description.

\begin{promptbox}[title=\textbf{Terminal-Bench 2.0 Task Instruction}]
Improve the pass rate of a terminal agent on Terminal-Bench~2.0.
The current baseline achieves 58.33\% on the 36-task development set
using GPT-5.5.

\smallskip
\noindent\textbf{Goal.}~Maximize pass rate on the development split.
Iterate on the development split exclusively.
Evaluate the held-out test split only after achieving a meaningful
improvement on development.

\smallskip
\noindent\textbf{Evaluation.}~Development: \texttt{HARBOR\_N\_CONCURRENT=8 python3 run\_eval.py --data data/dev.json}.
Test: \texttt{HARBOR\_N\_CONCURRENT=8 python3 run\_eval.py --data data/test.json}.
The score is the fraction of tasks solved correctly, with higher being better.

\smallskip
\noindent\textbf{Constraints.}~Modifiable components include the system
prompt, per-task extra instructions, the agent subclass, the base agent
and its ReAct loop, response parsers, terminal session management, and
new files under the agent or prompts directories.
The evaluation harness, task data files, API configuration, and baseline
reference record must not be modified.
\end{promptbox}

\subsection{BrowseComp}
\label{app:benchmark-browsecomp}

BrowseComp~\citep{browsecomp} is a search-agent benchmark for
multi-step question answering.
The material to be optimized is not the benchmark data itself, but a minimal
ReAct-style browsing harness~\citep{react} that answers BrowseComp questions
using search and page-reading tools.
The task measures whether an autonomous research agent can improve the control
logic around a search agent while leaving the evaluator and question sets
fixed.

\paragraph{Baseline.}
The initial material is a simple ReAct-based search harness centered
on \texttt{single\_agent\_gpt.py}.
For each question, a \texttt{GPTAgent} runs a single ReAct trajectory with a
system prompt, issues web-search and page-visit tool calls, and returns a
short final answer.
The online BrowseComp setting uses \texttt{SearchTool} and \texttt{VisitTool};
the local-corpus tool variants are available in the codebase but are not the
primary path for this benchmark.
The baseline uses the simple-evals BrowseComp query template and grader
template, with the same model family used for the answering agent and the
grader in the reported runs.

\paragraph{Development and test split.}
The development split contains 50 BrowseComp questions and is used for all
iterative optimization.
The held-out test split contains 300 non-overlapping BrowseComp questions and
is reserved for final verification.
In the released harness, these splits are exposed as
\texttt{data/bc\_val.jsonl} and \texttt{data/bc\_test.jsonl}.
Agents may inspect and evaluate on the development split during the search
loop, but they may not edit either data file or use the test split for
iteration.

\paragraph{Evaluator.}
The development evaluator is
\texttt{uv run python run\_eval.py --data data/bc\_val.jsonl --workers 8 --run-name \{node\_id\}}.
The held-out evaluator uses the same entry point with
\texttt{data/bc\_test.jsonl} and a distinct run name.
The evaluator sends each question to the harness, collects the final answer,
normalizes it through the fixed BrowseComp answer-grading path, and reports
accuracy as \texttt{correct}/\texttt{total}.
Items with repeated execution errors or unparsable final answers are counted as
errors and do not contribute to the correct count.
The primary metric is accuracy, with higher values better.

\paragraph{Search and tool configuration.}
The harness exposes two browsing tools to the answering agent.
\texttt{SearchTool} accepts one or more search queries and returns compact web
search results, while \texttt{VisitTool} fetches and cleans page content for a
selected URL.
The tool names and argument schemas are fixed so that alternative harness
designs remain compatible with the same ReAct agent interface.
Agents may change the system prompt, rollout orchestration, response parsing,
candidate aggregation, and other harness code, but the evaluation script,
question files, grader template, and reference answers are fixed.

\paragraph{Transfer-test protocol.}
After optimization on BrowseComp, the resulting search harness is frozen and
evaluated directly on unseen search-agent tasks, including HLE and
DeepSearchQA, without additional task-specific optimization.
The same code path, search/visit tool interface, and answer-production protocol
are reused.
This transfer test checks whether the discovered harness changes improve
general browsing behavior rather than merely fitting the BrowseComp development
questions.

\paragraph{Agent instruction.}
The agent is initialized with the following task description.

\begin{promptbox}[title=\textbf{BrowseComp Task Instruction}]
Optimize BrowseComp search-agent accuracy starting from the simple baseline on
the current main branch.
Do not use prior optimized branches, cached scores, or old result directories
as the starting point.

\smallskip
\noindent\textbf{Goal.}~Maximize answer accuracy on the BrowseComp
development split.
Use development feedback to propose and evaluate harness changes, and reserve
the held-out test split for final verification.

\smallskip
\noindent\textbf{Evaluation.}~Development:
\texttt{uv run python run\_eval.py --data data/bc\_val.jsonl --workers 8 --run-name \{node\_id\}}.
Test:
\texttt{uv run python run\_eval.py --data data/bc\_test.jsonl --workers 8 --run-name \{node\_id\}\_test}.
The evaluator uses the fixed simple-evals BrowseComp prompt and LLM grader.
The score is accuracy, with higher values better.

\smallskip
\noindent\textbf{Constraints.}~The agent may modify the search harness,
including \texttt{single\_agent\_gpt.py}, prompts, agent-control logic,
response parsers, rollout strategies, and aggregation code.
Do not modify \texttt{run\_eval.py}, the development or test data files, API
configuration, the grader template, or cached reference records.
\end{promptbox}

\subsection{Search-Agent Data Synthesis}
\label{app:benchmark-search-agent-data}

Search-Agent Data Synthesis is a pipeline-optimization benchmark that
evaluates whether a research agent can improve a synthetic question-answering
generation system.
The system begins from Wikipedia-derived topic seeds and produces
BrowseComp-style information-retrieval questions.
A high-quality generated item must satisfy three properties simultaneously:
structural well-formedness, answerability by a search-enabled model, and
sufficient difficulty that the solver does not answer it reliably in a single
attempt.
The task therefore measures not only whether generated questions are
answerable, but also whether they require multi-step evidence gathering rather
than surface-level lookup.

\paragraph{Baseline.}
The initial material is the unmodified default data-synthesis pipeline
distributed with the benchmark.
The baseline pipeline consists of five stages: extracting candidate facts from
a topic seed, rating facts by information value, constructing a question whose
answer is one selected fact, obfuscating the topic with supporting constraints,
and applying structural verification.
The generator, solver, and judge use GPT-5.5 through an OpenAI-compatible
interface.

\paragraph{Development and test split.}
The benchmark provides frozen seed snapshots for both evaluation splits.
The development split contains 50 seeds and is used exclusively for iterative
optimization.
The held-out test split contains 100 seeds and is reserved for final
verification after the search loop completes.
The pipeline generates at most one QA item per seed.
Fixing the seed snapshots makes the benchmark self-contained and ensures that
measured improvements reflect changes to the synthesis mechanism rather than
changes to the data distribution.

\paragraph{Evaluator.}
The evaluator is the benchmark's fixed \texttt{run\_eval.py} entry point.
It loads the requested seed split, runs the synthesis pipeline, solves each
generated item four times with a web-search-enabled ReAct model, and reports
a JSON summary on the final line of standard output.
The primary metric is:
\[
\mathrm{score}
\;=\;
\frac{1}{N}\sum_{i=1}^{N}
\bigl(\mathrm{pass@4} - \mathrm{pass@1}\bigr),
\]
where $\mathrm{pass@1}$ is $1$ if the first attempt is judged correct, and
$\mathrm{pass@4}$ is $1$ if at least one of four attempts is correct.
Higher scores indicate questions that remain solvable in principle but are
not answered reliably on the first attempt.

\paragraph{Agent instruction.}
The agent is initialized with the following task description.

\begin{promptbox}[title=\textbf{Search-Agent Data Synthesis Task Instruction}]
Improve the quality of generated search-intensive QA items by optimizing
the data-synthesis pipeline.
The goal metric rewards items that are solvable under repeated attempts but
not answered reliably on the first attempt:
\[
\mathrm{score}
\;=\;
\mathrm{mean}(\mathrm{pass@4} - \mathrm{pass@1}).
\]

\smallskip
\noindent\textbf{Goal.}~Maximize the score on the development split.
Prefer mechanism-level improvements that preserve well-formedness and
solvability while increasing multi-step evidence gathering.

\smallskip
\noindent\textbf{Evaluation.}~Development: \texttt{python run\_eval.py --split dev}.
Test: \texttt{python run\_eval.py --split test}.
Iterate on the development split. Run the test split only after a
meaningful improvement.

\smallskip
\noindent\textbf{Constraints.}~Only the synthesis pipeline under
\texttt{pipeline/} may be modified, including pipeline steps, the async
runner, the step registry, pipeline configuration, and prompt files.
The scoring module, evaluation harness, benchmark configuration, and
frozen seed snapshots must not be modified.
\end{promptbox}

\subsection{Math-Reasoning Data Synthesis}
\label{app:benchmark-math-reasoning-data}
\label{app:math-case-study}

Math-Reasoning Data Synthesis is a pipeline-optimization benchmark that
evaluates whether a research agent can improve a system for generating
mathematical contest problems.
The pipeline starts from fixed AIME/AIMO/NuminaMath-style seeds and produces
original problems with exact integer answers in the range 0--999.
The research target is the data-construction mechanism itself: generated
problems must be valid, novel, diverse, mathematically well specified, and
calibrated so that a fixed reference solver does not answer them trivially.

\paragraph{Baseline.}
The initial material is the unmodified default synthesis pipeline distributed
with the benchmark.
The pipeline takes each seed, generates several candidate contest problems,
filters candidates for answer format, rationale consistency, held-out overlap,
and near-duplicate overlap, and evaluates surviving candidates with a fixed
solver.
The default generator uses GPT-5.5-mini and the reference solver uses GPT-5.5.
The answer policy is \texttt{integer\_0\_999}.

\paragraph{Development and test split.}
The benchmark uses locked seed files for both evaluation splits.
The development split contains 10 seeds with 5 generated candidates per seed,
yielding up to 50 candidates total.
The held-out test split contains 12 seeds with 8 generated candidates per seed,
yielding up to 96 candidates total.
Each seed specifies a mathematical topic, technique, and target difficulty.
Agents iterate exclusively on the development split during the research loop.
the test split is reserved for milestone and final verification.

\paragraph{Evaluator.}
The evaluator is the fixed \texttt{run\_eval.py} entry point.
It loads the requested split, runs the synthesis pipeline, applies validity
and novelty filters, evaluates each surviving candidate with the fixed solver,
and prints a JSON metrics object on the final line of standard output.
The primary metric is:
\[
\mathrm{score}
\;=\;
\frac{1}{N}
\sum_{i=1}^{N}
\bigl(\mathrm{pass@4} - \mathrm{pass@1}\bigr),
\]
where $N$ is the number of generated candidates before filtering,
$\mathrm{pass@1}$ indicates whether the first solver sample is correct,
and $\mathrm{pass@4}$ indicates whether at least one of four solver
samples is correct.
Filtered candidates are treated as zero.
This metric rewards problems that are solvable under repeated attempts but
not immediately solved on the first sample.

\paragraph{Agent instruction.}
The agent is initialized with the following task description.

\begin{promptbox}[title=\textbf{Math-Reasoning Data Synthesis Task Instruction}]
Improve a pipeline that generates AIME/NuminaMath-style contest problems
with exact integer answers in the range 0--999.
Generated problems must be valid, novel, diverse, and calibrated so that the
fixed reference solver fails on the first sample but succeeds within four
samples.
The primary metric is $\mathrm{mean}(\mathrm{pass@4} - \mathrm{pass@1})$ over all
generated candidates.

\smallskip
\noindent\textbf{Goal.}~Maximize the primary metric reported on the final
standard-output line of the evaluator.
Prefer mechanism-level improvements such as structured or parametric
generation, programmatic answer computation, and difficulty calibration.

\smallskip
\noindent\textbf{Evaluation.}~Development: \texttt{uv run python run\_eval.py --split dev}.
Test: \texttt{uv run python run\_eval.py --split test}.
Iterate on the development split. Run the test split only for milestone or
final verification.

\smallskip
\noindent\textbf{Constraints.}~Modifiable files include
\texttt{configs/pipeline.yaml}, \texttt{prompts/generate\_problem.md},
\texttt{src/math\_synth\_bench/baseline.py}, and new modules under
\texttt{src/math\_synth\_bench/}.
The benchmark harness, seed files, evaluation references, metrics module,
and verification module must not be modified.
\end{promptbox}

\FloatBarrier
